%% 自動生成（tools/mknext.py）。各章の道しるべ「この章の先に」から、後ろの巻の章を集めたもの
\begin{longtblr}{colspec={Q[l,m,wd=58mm] X[l,m]}, width=\linewidth, rowhead=1, row{1}={bg=bkVol!14, font=\sffamily\bfseries\footnotesize},
  hlines={0.4pt,bkLine}, rowsep=1.5pt, cells={font=\footnotesize}}
この巻の章 & この道具を使う章 \\
{\bknextpill{6-2}\ 表計算で解こう：オイラー法} & {\bknextpill{7-10}\ 行列のくり返し：漸化式とオイラー法を行列で} \\
{\bknextpill{6-3}\ 積分するだけで解ける形} & {\bknextpill{8-13}\ 熱と波の伝わり方：偏微分方程式} \\
{\bknextpill{6-4}\ 自分に比例してへる：冷めるお湯} & {\bknextpill{9-5}\ 微分をかけ算に変える翻訳機：ラプラス変換\\ \bknextpill{9-9}\ 極と安定性：s の値でふるまいがわかる} \\
{\bknextpill{6-6}\ ばねと振り子} & {\bknextpill{7-8}\ 固有値と固有ベクトル：向きの変わらない方向} \\
{\bknextpill{6-8}\ 2乗してマイナスになる数：虚数} & {\bknextpill{9-9}\ 極と安定性：s の値でふるまいがわかる} \\
{\bknextpill{6-9}\ 複素平面：かけ算は回転と拡大} & {\bknextpill{7-6}\ 回転行列：くるっと回す} \\
{\bknextpill{6-10}\ オイラーの公式} & {\bknextpill{9-1}\ 波をたし合わせて形を作る：フーリエ級数\\ \bknextpill{9-3}\ 複素数で書くフーリエ級数} \\
{\bknextpill{6-11}\ 双曲線関数 cosh・sinh} & {\bknextpill{8-13}\ 熱と波の伝わり方：偏微分方程式} \\
{\bknextpill{6-12}\ 減衰振動を式で解く：特性方程式} & {\bknextpill{9-7}\ ラプラス変換で微分方程式を解く\\ \bknextpill{9-9}\ 極と安定性：s の値でふるまいがわかる} \\
{\bknextpill{6-13}\ ゆらし続けると：強制振動と共振} & {\bknextpill{9-6}\ 変換表と性質\\ \bknextpill{9-7}\ ラプラス変換で微分方程式を解く\\ \bknextpill{9-8}\ 伝達関数：入れると出てくるを s で表す\\ \bknextpill{9-12}\ フーリエとラプラスの使いみち} \\
{\bknextpill{6-14}\ 微分方程式がいくつも：連立微分方程式} & {\bknextpill{7-8}\ 固有値と固有ベクトル：向きの変わらない方向\\ \bknextpill{7-9}\ 対角化：まっすぐな方向に分けて考える\\ \bknextpill{7-10}\ 行列のくり返し：漸化式とオイラー法を行列で\\ \bknextpill{9-10}\ 連立微分方程式を行列で：状態方程式\\ \bknextpill{9-11}\ 行列の指数関数と固有値} \\
{\bknextpill{6-15}\ まとめ：物理は「微分方程式＋はじめの状態」で書ける} & {\bknextpill{7-0}\ 表のままで計算できたら？\\ \bknextpill{8-13}\ 熱と波の伝わり方：偏微分方程式\\ \bknextpill{9-5}\ 微分をかけ算に変える翻訳機：ラプラス変換\\ \bknextpill{9-10}\ 連立微分方程式を行列で：状態方程式} \\
\end{longtblr}
